\documentclass[10pt,a4paper]{amsart}
\usepackage[margin=19mm]{geometry}
\usepackage{amsmath,amssymb,mathtools}
\usepackage[scr=rsfs,cal=euler]{mathalfa}
\usepackage{xfrac}
\usepackage[unicode,hidelinks,bookmarks=false]{hyperref}
\newcommand{\ie}{i.e.\ }
\newcommand{\cf}{cf.\ }
\newcommand{\resp}{respectively\ }
\newcommand{\Subsection}{\subsection}
\providecommand{\nobreakdash}{\nobreak\hbox{-}}
\setlength{\emergencystretch}{2em}
\multlinegap0pt
\usepackage{bm}
\usepackage{bbm}
\usepackage{dsfont}
\usepackage{xspace}
\usepackage{cancel}
\usepackage{mathtools}
\usepackage{amsthm}
\makeatletter
\@ifundefined{theorem}{\newtheorem{theorem}{Theorem}}{}
\makeatother
\title[Higher Order Dualities between Prime Ideals]{Higher Order Dualities between Prime Ideals}
\author{Sroyon Sengupta}
\date{Source: arXiv:2512.22346v2 (CC BY 4.0)}
\begin{document}
\maketitle

\noindent\emph{Source and licence.} Higher Order Dualities between Prime Ideals. Version arXiv:2512.22346v2 (CC BY 4.0), distributed under the Creative Commons Attribution 4.0 International licence, as recorded in the arXiv licence metadata for this exact version and shown on the abstract page https://arxiv.org/abs/2512.22346v2. Full rights notice: licenses/arxiv-2512.22346-CC-BY-4.0.txt.\par\smallskip

\noindent\textit{Editorial scope.} Counted unit: the Theorem whose source label is \texttt{None} in the article (source0.tex lines 634--639, source label \texttt{None}), delivered together with the article's own complete original proof of it (source0.tex lines 640--783, one \texttt{proof} environment of 144 lines). No mathematics is written, repaired or re-derived: statement and proof are the article's own text line for line. Required context retained uncounted: none. Every definition, display and label of those lines is retained unchanged. Omitted from this package and not redistributed: 1--633, 784--1092. Nothing inside a retained excerpt is omitted or altered: every retained body is delivered exactly as the article's lines are written, so no omission receipt of any kind accompanies this package. Citations: the retained excerpts carry no citation command at all, so no bibliography entry of the article is transcribed and no reference is redistributed. Reference adaptation: none was needed and none was made; every reference inside the delivered unit resolves inside it, so no unresolved reference or citation is produced. Printed slips kept literal: no printed word, symbol, hypothesis or number of the counted statement or of its proof was altered. Layout and class: the article's own class is \texttt{article} with packages including geometry, hyperref, amssymb, amsmath, amsthm, mathrsfs, graphicx; the wrapper is the lane's pinned \texttt{amsart} editorial wrapper at 10pt on A4, which restates the theorem-like environments the retained text needs and adds the article's own macro definitions that the retained text needs (none), with extra wrapper packages (bm, bbm, dsfont, xspace, cancel, mathtools). The delivered numbering is the unit's own. Assets: no retained span contains a figure environment, a graphics-inclusion command, a PDF-page inclusion, a table or a TikZ picture; no image file is copied into this package, the package declares no asset dependency and no image file is redistributed with it. Licence: version arXiv:2512.22346v2 of this article is distributed under the Creative Commons Attribution 4.0 International licence, as recorded in the arXiv licence metadata for this exact version and shown on the abstract page https://arxiv.org/abs/2512.22346v2. A full rights notice is delivered at licenses/arxiv-2512.22346-CC-BY-4.0.txt. Characters: every retained excerpt is pure ASCII, so the delivered engine, XeLaTeX with its default Latin Modern text font, reports no missing character.
\begin{theorem}
    Let $L/K$ be a finite Galois extension with $G=Gal(L/K)$ and let $C \subset G$ be a fixed conjugacy class. Then for $X\geq 2$, we have
    \begin{align}
         \sum_{2 \leq N(I)\leq X} Q^2_C(I) = c_K\cdot \frac{|C|}{|G|}\cdot X + O\left(\frac{X(\log\log X)^2}{\log X}\right). \nonumber
    \end{align}
    \end{theorem}

    \begin{proof}
        Let us first recall the definition of $Q^2_C(I)$.
        \begin{align}
             Q_C^2(I) &:= \# \left\{\mathfrak{p} \subset \mathcal{O}_K :\; \mathfrak{p}\text{ is a prime; } I\subset \mathfrak{p}; \;M_2(I)=N(\mathfrak{p})\;\text{and}\;\left[\frac{L/K}{\mathfrak{p}}\right]=C,\;\mathfrak{p}\;\text{is unramified}\right\} .\nonumber
        \end{align}
        By definition, it is clear that ideals in $\mathcal{O}_K$ contained in prime ideals only of a particular norm do not contribute to the sum in the LHS. As done in the proof of \textit{Lemma 3.4}, in a similar fashion, we define the set $S_2(X,\mathfrak{p})$ as the set of ideals $I$ in $\mathcal{O}_K$ such that $N(I) \leq X$ with a unique and non-repeating prime ideal of the largest norm and $M_2(I)=N(\mathfrak{p})$. We note here that for any ideal $I \subset \mathfrak{p}$, for some prime $\mathfrak{p}$ such that $N(I)\leq X$ and $M_2(I)=N(\mathfrak{p})$, $N(\mathfrak{p}) \leq \sqrt{X}$. Therefore, using \textit{Lemma 3.1}, we have that
        \begin{align}
             \sum_{2\leq N(I)\leq X} Q_C^2(I) = \sum_{\substack{N(\mathfrak{p})\leq \sqrt{X} \\ \left[\frac{L/K}{\mathfrak{p}}\right]=C}} |S_2(X,\mathfrak{p})| +  O\left(\frac{X}{e^{c_3\sqrt{\log X\log\log X}}}\right).
        \end{align}
        We therefore move on to estimate the sum on the RHS of (4.1). We choose $I \in S_2(X,\mathfrak{p})$ such that $I=\mathfrak{M.pq}$, where $M_1(I)=\mathfrak{q}$, $M_2(I)=\mathfrak{p}$ and $\mathfrak{M} \subset \mathcal{O}_K$ is an ideal such that $M_1(\mathfrak{M})\leq N(\mathfrak{p})$. Thus, we write
        \begin{align}
            |S_2(X,\mathfrak{p})| = \sum_{\substack{N(\mathfrak{M}) \leq \frac{X}{N(\mathfrak{p})^2} \\ M_1(\mathfrak{M}) \leq N(\mathfrak{p})}} \sum_{\substack{N(\mathfrak{p})<N(\mathfrak{q}) \\ N(\mathfrak{M.pq}) \leq X}} 1 = \sum_{N(\mathfrak{p})<N(\mathfrak{q})\leq \frac{X}{N(\mathfrak{p})}} \sum_{\substack{N(\mathfrak{M})\leq \frac{X}{N(\mathfrak{pq})} \\ M_1(\mathfrak{M}) \leq N(\mathfrak{p})}} 1 .
        \end{align}
    Using the definition of $\Psi_K(X,Y)$, we then have that 
    \begin{align}
        |S_2(X,\mathfrak{p})| = \sum_{\substack{\mathfrak{q} \in \mathcal{O}_K \\ \mathfrak{q} \text{ is a prime} \\N(\mathfrak{p})<N(\mathfrak{q})\leq \frac{X}{N(\mathfrak{p})}}} \Psi_K\left(\frac{X}{N(\mathfrak{pq})},N(\mathfrak{p})\right).
    \end{align} 
    Now, we observe that for $Y \leq e^{(\log X)^{1-\delta}}$, for some $\delta  >0$, we have using (3.27) that
    \begin{align}
        \sum_{\substack{N(\mathfrak{p}) \leq Y \\ \left[\frac{L/K}{\mathfrak{p}}\right] =C}} |S_2(X,\mathfrak{p})| \leq \sum_{N(\mathfrak{p}) \leq Y} |S_2(X,\mathfrak{p})| \leq   \sum_{N(\mathfrak{p})\leq Y} \Psi_{K,2}\left({X},\mathfrak{p}\right) \ll \frac{X\log Y}{\log X}.
    \end{align}
        Therefore, we are left to estimate the sum of $|S_2(X,\mathfrak{p})|$ over the interval $Y<N(\mathfrak{p})\leq \sqrt{X}$. Therefore, from (4.3) we have
        \begin{align}
            \sum_{\substack{Y<N(\mathfrak{p}) \leq \sqrt{X} \\ \left[\frac{L/K}{\mathfrak{p}}\right] =C}} |S_2(X,\mathfrak{p})| = \sum_{\substack{Y<N(\mathfrak{p}) \leq \sqrt{X} \\ \left[\frac{L/K}{\mathfrak{p}}\right] =C}}\; \sum_{N(\mathfrak{p})<N(\mathfrak{q})\leq \frac{X}{N(\mathfrak{p})}} \Psi_K\left(\frac{X}{N(\mathfrak{p}\mathfrak{q})},N(\mathfrak{p})\right).
        \end{align}
        It is worth mentioning here that the idea of this proof is quite similar to the proof of \textit{Lemma 3.4}. We change the order of summation in the RHS of (4.5), which would require a split in the summation. We keep in mind the following conditions that $N(\mathfrak{p})$ need to satisfy:
        \begin{align}
            Y<N(\mathfrak{p}), \quad N(\mathfrak{p})<N(\mathfrak{q}), \quad N(\mathfrak{p})\leq \frac{X}{N(\mathfrak{q})},\quad N(\mathfrak{p}) \leq \sqrt{X} .\nonumber
        \end{align}
        Continuing from (4.5), we then get 
        \begin{align}
            \sum_{\substack{Y<N(\mathfrak{p}) \leq \sqrt{X} \\ \left[\frac{L/K}{\mathfrak{p}}\right] =C}} |S_2(X,\mathfrak{p})| &= \sum_{\substack{Y <N(\mathfrak{q}) \leq \sqrt{X}}} \;  \sum_{\substack{Y<N(\mathfrak{p})<N(\mathfrak{q}) \\ \left[\frac{L/K}{\mathfrak{p}}\right] =C}} \Psi_K\left(\frac{X}{N(\mathfrak{p}\mathfrak{q})},N(\mathfrak{p})\right) \nonumber \\ &\hspace{2cm}+ \sum_{\sqrt{X} < N(\mathfrak{q}) \leq \frac{X}{Y}} \;\sum_{\substack{Y<N(\mathfrak{p})\leq \frac{X}{N(\mathfrak{q})} \\ \left[\frac{L/K}{\mathfrak{p}}\right] =C} } \Psi_K\left(\frac{X}{N(\mathfrak{p}\mathfrak{q})},N(\mathfrak{p})\right) .
        \end{align}
        To bring in the Chebotarev Density, we now interchange the inner sums of both the double sums in the RHS of (4.6) with suitable integrals. This is indeed one of the most crucial steps of the whole proof. We do the replacements as follows:
        \begin{align}
            \sum_{\substack{Y<N(\mathfrak{p})<N(\mathfrak{q}) \\ \left[\frac{L/K}{\mathfrak{p}}\right] =C}} \Psi_K\left(\frac{X}{N(\mathfrak{p}\mathfrak{q})},N(\mathfrak{p})\right) \quad &\text{by} \quad \frac{|C|}{|G|}\int_Y^{N(\mathfrak{q})} \Psi_K\left(\frac{X}{tN(\mathfrak{q})},t\right)\frac{dt}{\log t} ,\nonumber \\
     \sum_{\substack{Y<N(\mathfrak{p})<\frac{X}{N(\mathfrak{q})} \\ \left[\frac{L/K}{\mathfrak{p}}\right] =C}} \Psi_K\left(\frac{X}{N(\mathfrak{p}\mathfrak{q})},N(\mathfrak{p})\right) \quad &\text{by} \quad \frac{|C|}{|G|}\int_Y^{\frac{X}{N(\mathfrak{q})}} \Psi_K\left(\frac{X}{tN(\mathfrak{q})},t\right)\frac{dt}{\log t}. \nonumber
        \end{align}
         Of course, it is now time to estimate the following errors that arise due to such a replacement:
    \begin{align}
        E_1&:=\sum_{Y<N(\mathfrak{q})\leq \sqrt{X}}\left(\sum_{\substack{Y<N(\mathfrak{p})<N(\mathfrak{q}) \\ \left[\frac{L/K}{\mathfrak{p}}\right] =C}} \Psi_K\left(\frac{X}{N(\mathfrak{p}\mathfrak{q})},N(\mathfrak{p})\right) - \frac{|C|}{|G|}\int_Y^{N(\mathfrak{q})} \Psi_K\left(\frac{X}{tN(\mathfrak{q})},t\right)\frac{dt}{\log t}\right) ,\nonumber \\
     E_2 &:=\sum_{\sqrt{X}<N(\mathfrak{q})\leq \frac{X}{Y}}\left(\sum_{\substack{Y<N(\mathfrak{p})<\frac{X}{N(\mathfrak{q})} \\ \left[\frac{L/K}{\mathfrak{p}}\right] =C}} \Psi_K\left(\frac{X}{N(\mathfrak{p}\mathfrak{q})},N(\mathfrak{p})\right)- \frac{|C|}{|G|}\int_Y^{\frac{X}{N(\mathfrak{q})}} \Psi_K\left(\frac{X}{tN(\mathfrak{q})},t\right)\frac{dt}{\log t}\right) .\nonumber
    \end{align}
From (4.6), we then have 
\begin{align}
    \sum_{\substack{Y<N(\mathfrak{p}) \leq \sqrt{X} \\ \left[\frac{L/K}{\mathfrak{p}}\right] =C}} |S_2(X,\mathfrak{p})| = \frac{|C|}{|G|}\sum_{Y<N(\mathfrak{q})\leq \sqrt{X}}\int_Y^{N(\mathfrak{q})} \Psi_K\left(\frac{X}{tN(\mathfrak{q})},t\right)\frac{dt}{\log t}& \nonumber \\+
\frac{|C|}{|G|}\sum_{\sqrt{X}<N(\mathfrak{q})\leq \frac{X}{Y}}\int_Y^{\frac{X}{N(\mathfrak{q})}}& \Psi_K\left(\frac{X}{tN(\mathfrak{q})},t\right)\frac{dt}{\log t}\nonumber \\&+E_1+E_2 .
\end{align}
To estimate the error terms $E_1$ and $E_2$, we use the strong form of Chebotarev Density Theorem (\textit{Theorem 1.2}). Changing the order of summations in $E_1$, we observe that
\begin{align}
    |E_1| &= \left| \sum_{Y<N(\mathfrak{q})\leq \sqrt{X}}\left(\sum_{\substack{Y<N(\mathfrak{p})<N(\mathfrak{q}) \\ \left[\frac{L/K}{\mathfrak{p}}\right] =C}} \Psi_K\left(\frac{X}{N(\mathfrak{p}\mathfrak{q})},N(\mathfrak{p})\right) - \frac{|C|}{|G|}\int_Y^{N(\mathfrak{q})} \Psi_K\left(\frac{X}{tN(\mathfrak{q})},t\right)\frac{dt}{\log t}\right)\right|  \nonumber \\
    &= \left| \sum_{Y<N(\mathfrak{q})\leq \sqrt{X}}\left(\sum_{\substack{Y<N(\mathfrak{p})<N(\mathfrak{q}) \\ \left[\frac{L/K}{\mathfrak{p}}\right] =C}}\; \sum_{\substack{N(\mathfrak{M}) \leq \frac{X}{N(\mathfrak{pq})} \\ M_1(\mathfrak{M})\leq N(\mathfrak{p})}} 1- \frac{|C|}{|G|}\int_Y^{N(\mathfrak{q})} \left\{\sum_{\substack{N(\mathfrak{M}) \leq \frac{X}{tN(\mathfrak{q})} \\ M_1(\mathfrak{M})\leq t}}1\right\}\frac{dt}{\log t}\right)\right| \nonumber \\
    &=    \sum_{Y<N(\mathfrak{q})\leq \sqrt{X}} \;\sum_{N(\mathfrak{M}) \leq \frac{X}{YN(\mathfrak{q})}}\left| \sum_{\substack{\max(Y,M_1(\mathfrak{M}))\leq N(\mathfrak{p})\leq  \min\left(N(\mathfrak{q}),\frac{X}{N(\mathfrak{Mq})}\right) \\  \left[\frac{L/K}{\mathfrak{p}}\right] =C            }       }      1    - \frac{|C|}{|G|}\int_{\max(Y,M_1(\mathfrak{M}))}^{\min\left(N(\mathfrak{q}),\frac{X}{N(\mathfrak{Mq})}\right)}\frac{dt}{\log t}    \right|.
\end{align}
Of course, the innermost summand involving the absolute value resembles the LHS of \textit{Theorem 1.2}. Also, we note that error function $\frac{X}{e^{c_1\sqrt{\frac{\log X}{n_K}}}}$ is increasing, and therefore, we use the CDT with the bound $\frac{X}{N(\mathfrak{Mq})}$ on $N(\mathfrak{p})$. We denote the constant $c_4 = \frac{c_1}{\sqrt{n_K}}$. Thus, from (4.8), we get
\begin{align}
     |E_1| &\ll \sum_{Y<N(\mathfrak{q})\leq \sqrt{X}} \;\sum_{N(\mathfrak{M}) \leq \frac{X}{YN(\mathfrak{q})}} \frac{X}{N(\mathfrak{Mq})e^{c_4\sqrt{\log \left(\frac{X}{N(\mathfrak{Mq})}\right)}}} \nonumber \\
    &\ll \frac{X}{e^{c_4\sqrt{\log Y}}}\sum_{Y<N(\mathfrak{q})\leq \sqrt{X}}\frac{1}{N(\mathfrak{q})} \;\sum_{N(\mathfrak{M}) \leq \frac{X}{YN(\mathfrak{q})}}\frac{1}{N(\mathfrak{M})} \nonumber \\ 
    &\ll \frac{X\log X\log\log X}{e^{c_4\sqrt{\log Y}}}.
\end{align}
Similarly, for $E_2$, we have that
\begin{align}
    |E_2| &= \left| \sum_{\sqrt{X}<N(\mathfrak{q})\leq\frac{X}{Y}}\left(\sum_{\substack{Y<N(\mathfrak{p})\leq \frac{X}{N(\mathfrak{q})} \\ \left[\frac{L/K}{\mathfrak{p}}\right] =C}} \Psi_K\left(\frac{X}{N(\mathfrak{p}\mathfrak{q})},N(\mathfrak{p})\right) - \frac{|C|}{|G|}\int_Y^{\frac{X}{N(\mathfrak{q})}} \Psi_K\left(\frac{X}{tN(\mathfrak{q})},t\right)\frac{dt}{\log t}\right)\right|  \nonumber \\
    &= \left| \sum_{Y<N(\mathfrak{q})\leq \frac{X}{Y}}\left(\sum_{\substack{Y<N(\mathfrak{p})\leq \frac{X}{N(\mathfrak{q})} \\ \left[\frac{L/K}{\mathfrak{p}}\right] =C}}\; \sum_{\substack{N(\mathfrak{M}) \leq \frac{X}{N(\mathfrak{pq})} \\ M_1(\mathfrak{M})\leq N(\mathfrak{p})}} 1- \frac{|C|}{|G|}\int_Y^{\frac{X}{N(\mathfrak{q})}} \left\{\sum_{\substack{N(\mathfrak{M}) \leq \frac{X}{tN(\mathfrak{q})} \\ M_1(\mathfrak{M})\leq t}}1\right\}\frac{dt}{\log t}\right)\right| \nonumber \\
    &= \sum_{Y<N(\mathfrak{q})\leq \frac{X}{Y}} \sum_{N(\mathfrak{M}) \leq \frac{X}{YN(\mathfrak{q})}}\left|           \sum_{\substack{\max(Y,M_1(\mathfrak{M}))\leq N(\mathfrak{p})\leq  \min\left(\frac{X}{N(\mathfrak{q})},\frac{X}{N(\mathfrak{Mq})}\right) \\  \left[\frac{L/K}{\mathfrak{p}}\right] =C            }       }      1    - \frac{|C|}{|G|}\int_{\max(Y,M_1(\mathfrak{M}))}^{\min\left(\frac{X}{N(\mathfrak{q})},\frac{X}{N(\mathfrak{Mq})}\right)}\frac{dt}{\log t}          \right| \nonumber \\
     &\ll \sum_{\sqrt{X}<N(\mathfrak{q}) \leq \frac{X}{Y}}\sum_{N(\mathfrak{M}) \leq \frac{X}{YN(\mathfrak{q})}} \frac{X}{N(\mathfrak{Mq})e^{c_4\sqrt{\log\left(\frac{X}{N(\mathfrak{Mq})}\right)}}} \nonumber \\
    &\ll \frac{X}{e^{c_4\sqrt{\log Y}}}\sum_{\sqrt{X}<N(\mathfrak{q})\leq \frac{X}{Y}}\frac{1}{N(\mathfrak{q})}\sum_{N(\mathfrak{M})\leq \frac{X}{YN(\mathfrak{q})}}\frac{1}{N(\mathfrak{M})} \nonumber \\
    &\ll \frac{X\log X\log \log X}{e^{c_4\sqrt{\log Y}}}.
\end{align}
Therefore, combining (4.7), (4.9), and (4.10), we have that
\begin{align}
     \sum_{\substack{Y<N(\mathfrak{p}) \leq \sqrt{X} \\ \left[\frac{L/K}{\mathfrak{p}}\right] =C}} |S_2(X,\mathfrak{p})| = \frac{|C|}{|G|}&\sum_{Y<N(\mathfrak{q})\leq \sqrt{X}}\int_Y^{N(\mathfrak{q})} \Psi_K\left(\frac{X}{tN(\mathfrak{q})},t\right)\frac{dt}{\log t} \nonumber \\+
&\frac{|C|}{|G|}\sum_{\sqrt{X}<N(\mathfrak{q})\leq \frac{X}{Y}}\int_Y^{\frac{X}{N(\mathfrak{q})}} \Psi_K\left(\frac{X}{tN(\mathfrak{q})},t\right)\frac{dt}{\log t}+O\left(\frac{X\log X\log \log X}{e^{c_4\sqrt{\log Y}}}\right).
\end{align}
Therefore, along with (4.4), we get from (4.11) that
\begin{align}
     \sum_{\substack{N(\mathfrak{p}) \leq \sqrt{X} \\ \left[\frac{L/K}{\mathfrak{p}}\right] =C}} |S_2(X,\mathfrak{p})| = \frac{|C|}{|G|}\left[\sum_{Y<N(\mathfrak{q})\leq \sqrt{X}}\right.&\left.\int_Y^{N(\mathfrak{q})} \Psi_K\left(\frac{X}{tN(\mathfrak{q})},t\right)\frac{dt}{\log t} +
\sum_{\sqrt{X}<N(\mathfrak{q})\leq \frac{X}{Y}}\int_Y^{\frac{X}{N(\mathfrak{q})}} \Psi_K\left(\frac{X}{tN(\mathfrak{q})},t\right)\frac{dt}{\log t}\right]\nonumber \\&+O\left(\frac{X\log Y}{\log X}\right)+O\left(\frac{X\log X\log \log X}{e^{c_4\sqrt{\log Y}}}\right).
\end{align}
Now, we will evaluate the pseudo-integro-sums using further estimates. Surprisingly, the very same idea used above will work here too. Let us recall the definition of $Q^2(I)$:
\begin{align}
    Q^2(I)&:=\# \{\mathfrak{p}\subset  \mathcal{O}_K:\;I \subset \mathfrak{p};\;M_2(I)=Nm(\mathfrak{p})\} . \nonumber
\end{align}
Note that by the definition of $S_2(X,\mathfrak{p})$, we can write using \textit{Lemma 3.1}, in a similar fashion as of (4.1) that
\begin{align}
    \sum_{2\leq N(I)\leq X} Q^2(I) = \sum_{N(\mathfrak{p})\leq \sqrt{X} } |S_2(X,\mathfrak{p})| +  O\left(\frac{X}{e^{c_3\sqrt{\log X\log\log X}}}\right).
\end{align}
Here, we observe that the sum on the right of (4.13) is exactly the same as that we treated in the first half of the proof, without the conjugacy class condition on the prime ideal $\mathfrak{p}$. But the change in the consideration of the condition hardly changes anything in the estimation of the sum. We can start by using the estimate in (3.27) to see that for $Y\leq e^{(\log X)^{1-\delta}}$, for some $\delta>0$, 
\begin{align}
    \sum_{N(\mathfrak{p})\leq Y} |S_2(X,\mathfrak{p})| \ll\frac{X\log Y}{\log X}.
\end{align}
Using the idea in (4.3), we can write here that
\begin{align}
    \sum_{Y<N(\mathfrak{p}) \leq \sqrt{X}} |S_2(X,\mathfrak{p})| = \sum_{\substack{Y<N(\mathfrak{p}) \leq \sqrt{X} }} \sum_{N(\mathfrak{p})<N(\mathfrak{q})\leq \frac{X}{N(\mathfrak{p})}} \Psi_K\left(\frac{X}{N(\mathfrak{p}\mathfrak{q})},N(\mathfrak{p})\right) .
\end{align}
As done earlier, we change the order of summation in (4.15) and then split the double sum into two to get
\begin{align}
    \sum_{\substack{Y<N(\mathfrak{p}) \leq \sqrt{X} }} |S_2(X,\mathfrak{p})| &= \sum_{\substack{Y <N(\mathfrak{q}) \leq \sqrt{X}}} \;  \sum_{\substack{Y<N(\mathfrak{p})<N(\mathfrak{q}) }} \Psi_K\left(\frac{X}{N(\mathfrak{p}\mathfrak{q})},N(\mathfrak{p})\right) \nonumber \\ &\hspace{2cm}+ \sum_{\sqrt{X} < N(\mathfrak{q}) \leq \frac{X}{Y}} \;\sum_{\substack{Y<N(\mathfrak{p})\leq \frac{X}{N(\mathfrak{q})}} } \Psi_K\left(\frac{X}{N(\mathfrak{p}\mathfrak{q})},N(\mathfrak{p})\right) .
\end{align}
We observe that the only difference between equations (4.6) and (4.16) is the conjugacy class condition. Thus, we do a similar replacement of the inner sums with integrals, but this time, without the Chebotarev Density factor in front of the integrals. We make the following replacements:
\begin{align}
    \sum_{\substack{Y<N(\mathfrak{p})<N(\mathfrak{q}) }} \Psi_K\left(\frac{X}{N(\mathfrak{p}\mathfrak{q})},N(\mathfrak{p})\right) \quad &\text{by} \quad \int_Y^{N(\mathfrak{q})} \Psi_K\left(\frac{X}{tN(\mathfrak{q})},t\right)\frac{dt}{\log t}, \nonumber \\
     \sum_{\substack{Y<N(\mathfrak{p})<\frac{X}{N(\mathfrak{q})} }} \Psi_K\left(\frac{X}{N(\mathfrak{p}\mathfrak{q})},N(\mathfrak{p})\right) \quad &\text{by} \quad \int_Y^{\frac{X}{N(\mathfrak{q})}} \Psi_K\left(\frac{X}{tN(\mathfrak{q})},t\right)\frac{dt}{\log t}. \nonumber
\end{align}
Of course, we then have the errors that need to be estimated. We denote them as:
\begin{align}
    E_3&:=\sum_{Y<N(\mathfrak{q})\leq \sqrt{X}}\left(\sum_{\substack{Y<N(\mathfrak{p})<N(\mathfrak{q}) }} \Psi_K\left(\frac{X}{N(\mathfrak{p}\mathfrak{q})},N(\mathfrak{p})\right) - \int_Y^{N(\mathfrak{q})} \Psi_K\left(\frac{X}{tN(\mathfrak{q})},t\right)\frac{dt}{\log t}\right), \nonumber \\
     E_4 &:=\sum_{\sqrt{X}<N(\mathfrak{q})\leq \frac{X}{Y}}\left(\sum_{\substack{Y<N(\mathfrak{p})<\frac{X}{N(\mathfrak{q})}}} \Psi_K\left(\frac{X}{N(\mathfrak{p}\mathfrak{q})},N(\mathfrak{p})\right)-\int_Y^{\frac{X}{N(\mathfrak{q})}} \Psi_K\left(\frac{X}{tN(\mathfrak{q})},t\right)\frac{dt}{\log t}\right). \nonumber
\end{align}
For the estimation of $E_3$ and $E_4$, we use \textit{Theorem 1.3} and other standard summation estimates, as used in the proof above, and also in the proof of \textit{Lemma 3.4}. Following similar modifications and calculations, we finally get that
\begin{align}
    |E_3| \ll \frac{X\log X \log\log X}{e^{c_1\sqrt{\log Y}}} \quad \text{and} \quad |E_4| \ll \frac{X\log X \log\log X}{e^{c_1\sqrt{\log Y}}}.
\end{align}
Therefore, combining (4.14), (4.16), and (4.17), we have that
\begin{align}
    \sum_{\substack{N(\mathfrak{p}) \leq \sqrt{X} }} |S_2(X,\mathfrak{p})| = \sum_{Y<N(\mathfrak{q})\leq \sqrt{X}}&\int_Y^{N(\mathfrak{q})} \Psi_K\left(\frac{X}{tN(\mathfrak{q})},t\right)\frac{dt}{\log t} +
\sum_{\sqrt{X}<N(\mathfrak{q})\leq \frac{X}{Y}}\int_Y^{\frac{X}{N(\mathfrak{q})}} \Psi_K\left(\frac{X}{tN(\mathfrak{q})},t\right)\frac{dt}{\log t}\nonumber \\&+O\left(\frac{X\log Y}{\log X}\right)+O\left(\frac{X\log X\log \log X}{e^{c_1\sqrt{\log Y}}}\right).
\end{align}
Further, using (4.13), we can write from (4.18) that
\begin{align}
    \sum_{2\leq N(I)\leq X} Q^2(I) = & \sum_{Y<N(\mathfrak{q})\leq \sqrt{X}}\int_Y^{N(\mathfrak{q})} \Psi_K\left(\frac{X}{tN(\mathfrak{q})},t\right)\frac{dt}{\log t} +
\sum_{\sqrt{X}<N(\mathfrak{q})\leq \frac{X}{Y}}\int_Y^{\frac{X}{N(\mathfrak{q})}} \Psi_K\left(\frac{X}{tN(\mathfrak{q})},t\right)\frac{dt}{\log t}\nonumber \\&+O\left(\frac{X\log Y}{\log X}\right)+O\left(\frac{X\log X\log \log X}{e^{c_1\sqrt{\log Y}}}\right)+  O\left(\frac{x}{e^{c_3\sqrt{\log X\log\log X}}}\right).
\end{align}
We can rewrite the LHS of (4.19), using \textit{Theorem 1.5}, as
\begin{align}
     \sum_{2\leq N(I)\leq X} Q^2(I) = \sum_{2 \leq N(I)\leq X}1+\sum_{\substack{2\leq N(I) \leq X \\ Q^2(I) \geq 2}} (Q^2(I)-1) = c_K\cdot X + \sum_{\substack{2\leq N(I) \leq X \\ Q^2(I) \geq 2}} (Q^2(I)-1)  + O\left(X^{1-\frac{1}{d}}\right),
\end{align}
where $d = [K:\mathbb{Q}]$. Then using \textit{Lemma 3.4}, we have that
\begin{align}
     \sum_{2\leq N(I)\leq X} Q^2(I) =
c_K\cdot X +O\left(\frac{X\log Y}{\log X}\right)+O\left(\frac{X\log\log X}{e^{c_1\sqrt{\log Y}}}\right) .
\end{align}
Thus, combining (4.19) and (4.21) and rearranging the terms, we get
\begin{align}
    &\sum_{Y<N(\mathfrak{q})\leq \sqrt{X}}\int_Y^{N(\mathfrak{q})} \Psi_K\left(\frac{X}{tN(\mathfrak{q})},t\right)\frac{dt}{\log t} +
\sum_{\sqrt{X}<N(\mathfrak{q})\leq \frac{X}{Y}}\int_Y^{\frac{X}{N(\mathfrak{q})}} \Psi_K\left(\frac{X}{tN(\mathfrak{q})},t\right)\frac{dt}{\log t} \nonumber \\ 
&= c_K\cdot X +O\left(\frac{X\log Y}{\log X}\right)+O\left(\frac{X\log X\log\log X}{e^{c_4\sqrt{\log Y}}}\right).
\end{align}
Interestingly, the LHS of (4.22) is exactly the expression in brackets in the RHS of (4.12) that we were required to estimate. Therefore, now combining (4.1), (4.12), and (4.22), we finally get that
\begin{align}
     \sum_{2\leq N(I)\leq X} Q_C^2(I) = c_K\cdot \frac{|C|}{|G|}\cdot X  +O\left(\frac{X\log Y}{\log X}\right)+O\left(\frac{X\log X\log\log X}{e^{c_4\sqrt{\log Y}}}\right).
\end{align}
  Thus, a suitable choice of $Y = e^{\left(\frac{2}{c_4}\log\log X\right)^2}$ gives our desired result.
    \end{proof}

\end{document}
